% =====================================================================
% CAPÍTULO 4: REQUERIMIENTOS
% =====================================================================
\chapter{Requerimientos del Sistema}
\label{ch:requerimientos}

\section{Actores del Sistema}

\begin{table}[H]
\centering
\caption{Definición de actores del \ac{SGDIC}}
\label{tab:actores}
\small
\begin{tabularx}{\textwidth}{@{}l X@{}}
\toprule
\rowcolor{azulClaro}
\textbf{Actor} & \textbf{Descripción y responsabilidades} \\
\midrule
Estudiante &
  Usuario final que solicita cupos, consulta horarios, envía y recibe mensajes,
  responde cuestionarios de horario y de materias, califica docentes, recibe
  notificaciones de faltas y presenta quejas o sugerencias. \\
Docente &
  Gestiona sus materias y horarios, registra faltas, se comunica con estudiantes
  y con la secretaría, recibe evaluaciones consolidadas. \\
Secretaría &
  Recibe y clasifica quejas y solicitudes, asigna responsables, coordina la
  comunicación con docentes, verifica resultados y mantiene la trazabilidad. \\
Departamento &
  Administra el catálogo de materias y cupos, configura periodos de evaluación,
  consulta analítica, verifica resultados y aprueba decisiones académicas. \\
Sistema de notificaciones &
  Actor secundario que despacha correos, mensajes internos y recordatorios. \\
Sistema central universitario &
  Actor secundario que provee la matrícula, catálogo oficial y autenticación. \\
\bottomrule
\end{tabularx}
\end{table}

\section{Requerimientos Funcionales}

\subsection{Módulo de autenticación y roles}

\begin{longtable}{@{}p{1.5cm} p{6.5cm} p{4.5cm}@{}}
\caption{Requerimientos funcionales --- autenticación y roles}
\label{tab:rf-auth}\\
\toprule
\rowcolor{azulClaro}
\textbf{ID} & \textbf{Requerimiento} & \textbf{Actor} \\
\midrule
\endfirsthead
\toprule
\rowcolor{azulClaro}
\textbf{ID} & \textbf{Requerimiento} & \textbf{Actor} \\
\midrule
\endhead
RF-01 & El sistema debe autenticar usuarios mediante correo institucional y contraseña o credencial institucional. & Todos \\
RF-02 & El sistema debe asignar permisos según el rol (\ac{RBAC}): estudiante, docente, secretaría, departamento. & Todos \\
RF-03 & El sistema debe permitir la recuperación de contraseña sin intervención de la secretaría. & Todos \\
RF-04 & El sistema debe registrar bitácora de accesos y operaciones críticas. & Departamento \\
\bottomrule
\end{longtable}

\subsection{Módulo de cupos}

\begin{longtable}{@{}p{1.5cm} p{6.5cm} p{4.5cm}@{}}
\caption{Requerimientos funcionales --- gestión de cupos}
\label{tab:rf-cupos}\\
\toprule
\rowcolor{verdeClaro}
\textbf{ID} & \textbf{Requerimiento} & \textbf{Actor} \\
\midrule
\endfirsthead
\toprule
\rowcolor{verdeClaro}
\textbf{ID} & \textbf{Requerimiento} & \textbf{Actor} \\
\midrule
\endhead
RF-05 & El estudiante debe poder solicitar cupo en una materia mostrando la disponibilidad en tiempo real. & Estudiante \\
RF-06 & El sistema debe validar prerrequisitos, créditos y conflictos de horario antes de aceptar la solicitud. & Sistema \\
RF-07 & El sistema debe aplicar una regla de priorización configurable (promedio, avance curricular, orden de solicitud). & Departamento \\
RF-08 & El sistema debe notificar al estudiante el resultado de la solicitud (aprobada, en lista de espera, rechazada). & Sistema \\
RF-09 & El sistema debe liberar automáticamente cupos de solicitudes no confirmadas en el plazo definido. & Sistema \\
RF-10 & La secretaría debe poder ver y gestionar la lista de espera por materia. & Secretaría \\
\bottomrule
\end{longtable}

\subsection{Módulo de materias}

\begin{longtable}{@{}p{1.5cm} p{6.5cm} p{4.5cm}@{}}
\caption{Requerimientos funcionales --- materias}
\label{tab:rf-materias}\\
\toprule
\rowcolor{naranjaClaro}
\textbf{ID} & \textbf{Requerimiento} & \textbf{Actor} \\
\midrule
\endfirsthead
\toprule
\rowcolor{naranjaClaro}
\textbf{ID} & \textbf{Requerimiento} & \textbf{Actor} \\
\midrule
\endhead
RF-11 & El docente debe poder proponer la creación de una materia nueva con su justificación. & Docente \\
RF-12 & El docente debe poder solicitar la edición de una materia existente (contenido, créditos, prerrequisitos, horario). & Docente \\
RF-13 & El sistema debe gestionar el flujo de aprobación: propuesta, revisión del departamento, aprobación o rechazo, publicación. & Departamento \\
RF-14 & El sistema debe mantener historial de versiones de cada materia con autor y fecha del cambio. & Sistema \\
RF-15 & El sistema debe publicar automáticamente las materias aprobadas en el catálogo visible para los estudiantes. & Sistema \\
\bottomrule
\end{longtable}

\subsection{Módulo de mensajería}

\begin{longtable}{@{}p{1.5cm} p{6.5cm} p{4.5cm}@{}}
\caption{Requerimientos funcionales --- mensajería}
\label{tab:rf-mensajes}\\
\toprule
\rowcolor{moradoClaro}
\textbf{ID} & \textbf{Requerimiento} & \textbf{Actor} \\
\midrule
\endfirsthead
\toprule
\rowcolor{moradoClaro}
\textbf{ID} & \textbf{Requerimiento} & \textbf{Actor} \\
\midrule
\endhead
RF-16 & El estudiante debe poder enviar mensajes a docentes de las materias en las que está inscrito. & Estudiante \\
RF-17 & El docente debe poder enviar mensajes masivos o individuales a los estudiantes de sus materias. & Docente \\
RF-18 & El sistema debe registrar estado del mensaje: enviado, entregado y leído. & Sistema \\
RF-19 & El sistema debe alertar cuando un mensaje no tenga respuesta en el plazo configurado. & Sistema \\
RF-20 & La secretaría debe poder intervenir en un hilo cuando el caso escale. & Secretaría \\
\bottomrule
\end{longtable}

\subsection{Módulo de faltas}

\begin{longtable}{@{}p{1.5cm} p{6.5cm} p{4.5cm}@{}}
\caption{Requerimientos funcionales --- registro de faltas}
\label{tab:rf-faltas}\\
\toprule
\rowcolor{azulClaro}
\textbf{ID} & \textbf{Requerimiento} & \textbf{Actor} \\
\midrule
\endfirsthead
\toprule
\rowcolor{azulClaro}
\textbf{ID} & \textbf{Requerimiento} & \textbf{Actor} \\
\midrule
\endhead
RF-21 & El docente debe poder registrar la asistencia por sesión de clase de forma rápida (lista o código). & Docente \\
RF-22 & El sistema debe notificar al estudiante el mismo día en que se registra su falta. & Sistema \\
RF-23 & El estudiante debe poder presentar una justificación adjuntando evidencia. & Estudiante \\
RF-24 & El docente debe poder aprobar o rechazar la justificación con comentario. & Docente \\
RF-25 & El sistema debe calcular el porcentaje de inasistencia acumulado por materia. & Sistema \\
RF-26 & El sistema debe alertar al estudiante, docente y secretaría cuando el acumulado supere el umbral permitido. & Sistema \\
\bottomrule
\end{longtable}

\subsection{Módulo de evaluación docente}

\begin{longtable}{@{}p{1.5cm} p{6.5cm} p{4.5cm}@{}}
\caption{Requerimientos funcionales --- evaluación docente}
\label{tab:rf-evaluacion}\\
\toprule
\rowcolor{verdeClaro}
\textbf{ID} & \textbf{Requerimiento} & \textbf{Actor} \\
\midrule
\endfirsthead
\toprule
\rowcolor{verdeClaro}
\textbf{ID} & \textbf{Requerimiento} & \textbf{Actor} \\
\midrule
\endhead
RF-27 & El departamento debe poder configurar el periodo, las preguntas y la fecha límite de la evaluación docente. & Departamento \\
RF-28 & El sistema debe habilitar la evaluación automáticamente al inicio del periodo configurado. & Sistema \\
RF-29 & El sistema debe enviar recordatorios automáticos a los estudiantes que no hayan completado la evaluación. & Sistema \\
RF-30 & El sistema debe garantizar el anonimato total del evaluador. & Sistema \\
RF-31 & El sistema debe bloquear servicios académicos hasta completar la evaluación, si así se configura. & Sistema \\
RF-32 & El sistema debe consolidar resultados y publicarlos solo al cumplirse la fecha límite o el umbral de participación. & Sistema \\
RF-33 & El docente debe poder consultar su resultado consolidado y su comparación histórica. & Docente \\
RF-34 & El departamento debe poder identificar docentes con resultados atípicos y generar planes de acompañamiento. & Departamento \\
\bottomrule
\end{longtable}

\subsection{Módulo de quejas y sugerencias}

\begin{longtable}{@{}p{1.5cm} p{6.5cm} p{4.5cm}@{}}
\caption{Requerimientos funcionales --- quejas y sugerencias}
\label{tab:rf-quejas}\\
\toprule
\rowcolor{naranjaClaro}
\textbf{ID} & \textbf{Requerimiento} & \textbf{Actor} \\
\midrule
\endfirsthead
\toprule
\rowcolor{naranjaClaro}
\textbf{ID} & \textbf{Requerimiento} & \textbf{Actor} \\
\midrule
\endhead
RF-35 & Cualquier actor debe poder registrar una queja o sugerencia con categoría y descripción. & Todos \\
RF-36 & El sistema debe clasificar automáticamente el caso por categoría, urgencia y área responsable. & Sistema \\
RF-37 & La secretaría debe poder asignar responsable, plazo y estado a cada caso. & Secretaría \\
RF-38 & La secretaría debe poder comunicarse formalmente con el docente involucrado dentro del caso. & Secretaría \\
RF-39 & El sistema debe notificar al solicitante cada cambio de estado del caso. & Sistema \\
RF-40 & El sistema debe permitir el cierre con evidencia de la acción tomada y encuesta de satisfacción. & Secretaría \\
RF-41 & El sistema debe permitir el anonimato del denunciante cuando la naturaleza del caso lo requiera. & Sistema \\
RF-42 & El sistema debe impedir el cierre de un caso sin acción registrada. & Sistema \\
\bottomrule
\end{longtable}

\subsection{Módulo de cuestionarios de materias y horarios}

\begin{longtable}{@{}p{1.5cm} p{6.5cm} p{4.5cm}@{}}
\caption{Requerimientos funcionales --- cuestionarios}
\label{tab:rf-cuestionarios}\\
\toprule
\rowcolor{moradoClaro}
\textbf{ID} & \textbf{Requerimiento} & \textbf{Actor} \\
\midrule
\endfirsthead
\toprule
\rowcolor{moradoClaro}
\textbf{ID} & \textbf{Requerimiento} & \textbf{Actor} \\
\midrule
\endhead
RF-43 & El sistema debe aplicar un cuestionario de disponibilidad horaria por estudiante y periodo. & Estudiante \\
RF-44 & El sistema debe aplicar un cuestionario de preferencia y evaluación de materias cursadas. & Estudiante \\
RF-45 & El sistema debe calcular la probabilidad de preferencia por cada franja horaria disponible. & Sistema \\
RF-46 & El sistema debe recomendar una asignación de horario que maximice la aceptación esperada. & Departamento \\
RF-47 & El sistema debe mostrar al departamento escenarios alternativos de horario con su impacto estimado. & Departamento \\
RF-48 & El sistema debe permitir al estudiante confirmar o rechazar el horario sugerido. & Estudiante \\
\bottomrule
\end{longtable}

\subsection{Módulo de analítica y soluciones rápidas}

\begin{longtable}{@{}p{1.5cm} p{6.5cm} p{4.5cm}@{}}
\caption{Requerimientos funcionales --- analítica}
\label{tab:rf-analitica}\\
\toprule
\rowcolor{azulClaro}
\textbf{ID} & \textbf{Requerimiento} & \textbf{Actor} \\
\midrule
\endfirsthead
\toprule
\rowcolor{azulClaro}
\textbf{ID} & \textbf{Requerimiento} & \textbf{Actor} \\
\midrule
\endhead
RF-49 & El sistema debe alimentar un repositorio analítico con los datos de todos los módulos. & Sistema \\
RF-50 & El sistema debe calcular los \ac{KPI} definidos en el Capítulo~\ref{ch:mision-vision}. & Sistema \\
RF-51 & El sistema debe detectar anomalías: cupos saturados, quejas recurrentes, docentes con baja evaluación, inasistencia atípica. & Sistema \\
RF-52 & El sistema debe generar recomendaciones accionables priorizadas por impacto y esfuerzo. & Sistema \\
RF-53 & El sistema debe estimar la demanda futura de cupos por materia. & Departamento \\
RF-54 & La secretaría y el departamento deben verificar resultados, quejas y evaluaciones en tableros consolidados. & Secretaría, Departamento \\
RF-55 & El sistema debe exportar reportes a \ac{CSV} y \ac{PDF}. & Secretaría, Departamento \\
RF-56 & El sistema debe registrar si la recomendación fue aplicada y medir su efecto posterior. & Sistema \\
\bottomrule
\end{longtable}

\section{Requerimientos No Funcionales}

\begin{table}[H]
\centering
\caption{Requerimientos no funcionales}
\label{tab:rnf}
\small
\begin{tabularx}{\textwidth}{@{}p{1.5cm} l X@{}}
\toprule
\rowcolor{verdeClaro}
\textbf{ID} & \textbf{Categoría} & \textbf{Descripción} \\
\midrule
RNF-01 & Usabilidad &
    Interfaz responsiva, accesible desde móvil, con un máximo de tres clics
    para completar cualquier operación frecuente. \\
RNF-02 & Rendimiento &
    Tiempo de respuesta menor a 2 segundos en el 95\% de las peticiones. \\
RNF-03 & Disponibilidad &
    Disponibilidad mínima del 99.5\% en horario académico (06:00--22:00). \\
RNF-04 & Seguridad &
    Cifrado en tránsito (TLS) y en reposo, autenticación multifactor para roles
    administrativos. \\
RNF-05 & Privacidad &
    Anonimización de evaluaciones y quejas sensibles; cumplimiento de normas de
    protección de datos personales. \\
RNF-06 & Escalabilidad &
    Soporte para el crecimiento de usuarios y datos sin rediseño estructural. \\
RNF-07 & Mantenibilidad &
    Arquitectura modular con separación de responsabilidades e interfaces
    documentadas. \\
RNF-08 & Interoperabilidad &
    \ac{API} \ac{REST} documentada para integrarse con el sistema central
    universitario. \\
RNF-09 & Trazabilidad &
    Registro de auditoría inmutable de cada transición de estado. \\
RNF-10 & Compatibilidad &
    Funcionamiento en navegadores modernos (Chrome, Edge, Firefox, Safari). \\
RNF-11 & Accesibilidad &
    Cumplimiento de pautas WCAG 2.1 nivel AA. \\
RNF-12 & Localización &
    Interfaz en español con soporte de zona horaria y formatos locales. \\
\bottomrule
\end{tabularx}
\end{table}

\section{Reglas de Negocio}

\begin{table}[H]
\centering
\caption{Reglas de negocio del \ac{DIC}}
\label{tab:reglas}
\small
\begin{tabularx}{\textwidth}{@{}p{1.3cm} l X@{}}
\toprule
\rowcolor{naranjaClaro}
\textbf{ID} & \textbf{Nombre} & \textbf{Regla} \\
\midrule
RN-01 & Prerrequisitos &
    No se puede otorgar cupo si el estudiante no ha aprobado los prerrequisitos
    declarados en la materia. \\
RN-02 & Cupo máximo &
    El número de cupos asignados nunca puede superar el cupo máximo definido
    para la materia y el periodo. \\
RN-03 & Conflicto horario &
    No se puede asignar cupo si el horario de la materia se solapa con otra
    materia ya inscrita por el estudiante. \\
RN-04 & Priorización &
    Ante escasez de cupos, se aplica prioridad por avance curricular, luego por
    promedio acumulado y finalmente por orden cronológico de la solicitud. \\
RN-05 & Anonimato de evaluación &
    El docente nunca accede a la identidad de quien lo evaluó; solo a resultados
    agregados con un mínimo de cinco respuestas. \\
RN-06 & Ventana de evaluación &
    La evaluación docente solo puede diligenciarse dentro del periodo
    configurado; fuera de él, el formulario se bloquea. \\
RN-07 & Notificación de falta &
    Toda falta registrada debe generar notificación al estudiante el mismo día. \\
RN-08 & Umbral de inasistencia &
    Al superar el 20\% de inasistencia en una materia, se genera alerta
    automática al estudiante, al docente y a la secretaría. \\
RN-09 & Cierre de queja &
    Una queja no puede cerrarse sin registrar la acción realizada y la fecha de
    cierre. \\
RN-10 & Escalamiento &
    Una queja crítica sin atender en 3 días se escala automáticamente a la
    dirección del departamento. \\
RN-11 & Edición de materia &
    Toda modificación curricular requiere revisión y aprobación del departamento
    antes de su publicación. \\
RN-12 & Cobertura de horario &
    Un horario solo se publica si alcanza el umbral mínimo de aceptación
    estimada definido por el departamento. \\
\bottomrule
\end{tabularx}
\end{table}

\section{Matriz de Trazabilidad Objetivos--Requerimientos}

\begin{table}[H]
\centering
\caption{Trazabilidad entre objetivos estratégicos y requerimientos}
\label{tab:trazabilidad}
\small
\begin{tabularx}{\textwidth}{@{}l X@{}}
\toprule
\rowcolor{moradoClaro}
\textbf{Objetivo} & \textbf{Requerimientos asociados} \\
\midrule
OE1 (Cupos) & RF-05, RF-06, RF-07, RF-08, RF-09, RF-10, RN-01, RN-02, RN-03, RN-04 \\
OE2 (Materias) & RF-11, RF-12, RF-13, RF-14, RF-15, RN-11 \\
OE3 (Comunicación) & RF-16, RF-17, RF-18, RF-19, RF-20 \\
OE4 (Faltas) & RF-21, RF-22, RF-23, RF-24, RF-25, RF-26, RN-07, RN-08 \\
OE5 (Evaluación) & RF-27 a RF-34, RN-05, RN-06 \\
OE6 (Quejas) & RF-35 a RF-42, RN-09, RN-10 \\
OE7 (Secretaría) & RF-10, RF-20, RF-37, RF-38, RF-40, RF-54, RF-55 \\
OE8 (Horarios) & RF-43, RF-44, RF-45, RF-46, RF-47, RF-48, RN-12 \\
OE9 (Verificación) & RF-50, RF-54, RF-55, RF-56 \\
OE10 (Analítica) & RF-49, RF-51, RF-52, RF-53, RF-56 \\
\bottomrule
\end{tabularx}
\end{table}